\section{Results}
\label{sec:results}

\begin{table}[tb]
\centering\small
\caption{Main results, all figures percentages, scored against the curated
reference. The first block compares with the published baseline on the
established split; the second gives the text-disjoint split, where the
held-out row was measured once after every decision was frozen.}
\label{tab:main}
\begin{tabular}{lrrrrrr}
\toprule
& \multicolumn{2}{c}{Scripted} & \multicolumn{2}{c}{Spontaneous} & \multicolumn{2}{c}{Mean} \\
\cmidrule(lr){2-3}\cmidrule(lr){4-5}\cmidrule(lr){6-7}
System & WER & CER & WER & CER & WER & CER \\
\midrule
\multicolumn{7}{l}{\emph{\texttt{sharedtask} development (speaker-disjoint)}} \\
Published baseline, converted & 16.08 & 3.77 & 10.23 & 1.39 & 13.15 & 2.58 \\
This work & 9.23 & 3.01 & 9.06 & 1.30 & \textbf{9.14} & 2.16 \\
\addlinespace
\multicolumn{7}{l}{\emph{\texttt{sharedtask\_both} (speaker- and text-disjoint), mean of three seeds}} \\
\quad development & 4.85 & 0.62 & 9.20 & 1.31 & 7.02 & 0.96 \\
\quad held-out & 10.79 & 2.33 & 3.38 & 0.39 & 7.08 & 1.36 \\
\bottomrule
\end{tabular}
\end{table}

\subsection{How long to train}
\label{sec:plateau}

Fine-tuning from an already-adapted checkpoint needs far less training than the
original 8{,}000 updates. On the text-disjoint split, every configuration
between 250 and 1{,}000 updates and between $1\times10^{-5}$ and
$5\times10^{-5}$ lands at 7.1--7.3 mean WER: below about 200 updates the
curated orthography is not yet learned, beyond about 1{,}000 the model
memorises, and by 8{,}000 it is worse than at 500. The useful conclusion is a
broad plateau, not an optimum.

Reaching that conclusion required a correction worth reporting. A single run at
500 updates and $2\times10^{-5}$ gave 6.96 mean WER, best of every
configuration on all four numbers, and we recorded it as the winner; a second
seed of the identical recipe gave 7.26. Three seeds measure the noise directly:
7.02 mean WER with standard deviation 0.19 on the text-disjoint development
set, and 9.23 with deviation 0.08 on the larger speaker-disjoint one, of which
Table~\ref{tab:main} reports the submitted seed. We treat any difference below
roughly 0.3 mean WER as indistinguishable from the random draw. The contrast between the two sets is easy to overlook when choosing one:
the smaller does not merely measure less precisely, it manufactures differences.

\subsection{Sentence overlap, measured}
\label{sec:overlap}

Section~\ref{sec:newsentences} makes the established split's text overlap a
measurement problem rather than a curiosity. The design holds the evaluation fixed and varies only what the model saw.
Condition~A trains on the text-disjoint split, so the development sentences
never appear. Condition~B adds the recordings of those same sentences made by
\emph{training} speakers, 3.65 hours more audio. Comparing A with B alone would
confound exposure with data volume, so condition~B$_{\text{eq}}$ removes 3.64
hours of randomly chosen A clips, leaving exactly A's 39.94 hours. Every
condition uses the same recipe, 250 updates, and the same development set.


Over two seeds each, A gives 5.20 and 4.97 scripted WER (0.65 and 0.62 CER)
and B$_{\text{eq}}$ gives 3.79 and 4.23 (0.48 and 0.55); B, with one seed and
more data, gives 4.68. The seed ranges of A and B$_{\text{eq}}$ do not overlap
-- the worse B$_{\text{eq}}$ run still beats the better A run -- so at equal
training volume, having heard a sentence read by other speakers is worth 1.1
WER points, a 21\% relative reduction, and about 19\% of CER.

That figure is a lower bound: it is measured at 250 updates, and memorisation
grows with training, so the original baseline's 8{,}000 updates should widen
it.

\subsection{Held-out evaluation}
\label{sec:heldout-results}


The lower block of Table~\ref{tab:main} gives the measurement. The means agree
closely, 7.02 against 7.08, which is the result we care about:
after dozens of decisions taken on the development set, the held-out estimate
did not collapse, and seed spread falls to 0.05.

The agreement hides two large opposing movements, both composition rather than
generalisation. Scripted error doubles because the held-out portion contains
one speaker who alone scores 98.67 WER over 59 clips; removing that same
speaker from the development set moves its scripted WER from 9.23 to 5.71.
Spontaneous error falls by two thirds because 46\% of the held-out spontaneous
portion comes from speakers already heard through silver transcriptions in
training. With 77 scripted speakers, a split places few voices in each
partition, so no such figure should be read without its per-speaker breakdown.

% =====================================================================
